\chapter{Numrat pseudorastësorë dhe kampionimi}
\begin{qellime}
Studenti shpjegon gjeneratorët pseudorastësorë, farën dhe gjendjen; zbaton transformimin e anasjelltë, Box--Muller dhe kampionimin me refuzim; teston shpërndarjen dhe pavarësinë; shmang përdorimin e gabuar paralel.
\end{qellime}

\section{Pseudo-rastësia}
Një gjenerator pseudorastësor është algoritëm determinist që prodhon një varg me veti statistike të ngjashme me mostrat e pavarura uniforme. Gjeneratori kongruencial linear është
\begin{equation}
 x_{n+1}=(ax_n+c)\bmod m,\qquad u_n=x_n/m.
\end{equation}
Zgjedhja e parametrave përcakton periodën dhe korrelacionet. Gjeneratorët modernë kanë gjendje më të madhe dhe teste më të mira, por nuk janë automatikisht kriptografikisht të sigurt.

Fara përcakton vargun. Ruajtja e saj mundëson riprodhim. Megjithatë, përdorimi i të njëjtës farë në procese paralele mund të krijojë vargje identike; duhen nënrryma të pavarura ose mekanizma \emph{spawn}.

\begin{figure}[ht]
\centering\includegraphics[width=.90\textwidth]{16_testi_cifteve_prng.pdf}
\caption{Grafiku $(u_n,u_{n+1})$ zbulon strukturë rrjete në një gjenerator të dobët.}
\end{figure}

\section{Testet statistike}
Histogrami kontrollon uniformitetin njëpërmasor, por nuk zbulon çdo varësi. Përdoren testi $\chi^2$, Kolmogorov--Smirnov, autokorrelacioni dhe teste spektrale. Një vlerë-$p$ nuk provon rastësinë; ajo vetëm mat përputhjen me një statistikë të zgjedhur. Testimi duhet të përfshijë shumë struktura dhe gjatësi vargu.

\section{Transformimi i anasjelltë}
Nëse $U\sim\mathcal U(0,1)$ dhe $F$ është kumulativi i synuar, atëherë
\begin{equation}
 X=F^{-1}(U)
\end{equation}
ka shpërndarje $F$. Për eksponencialen me normë $\lambda$,
\begin{equation}
 X=-\frac1\lambda\ln(1-U).
\end{equation}

\begin{figure}[ht]
\centering\includegraphics[width=.74\textwidth]{17_kampionimi_eksponencial.pdf}
\caption{Mostrat e gjeneruara me transformimin e anasjelltë dhe dendësia teorike.}
\end{figure}

\section{Mostrat Gauss}
Transformimi Box--Muller përdor dy uniforme të pavarura:
\begin{align}
 Z_1&=\sqrt{-2\ln U_1}\cos(2\pi U_2),\\
 Z_2&=\sqrt{-2\ln U_1}\sin(2\pi U_2).
\end{align}
Kjo rrjedh nga kalimi në koordinata polare të dendësisë Gauss dydimensionale. Bibliotekat moderne mund të përdorin algoritme më të shpejta, por parimi mbetet shembull i rëndësishëm i transformimit të ndryshoreve.

\section{Kampionimi me refuzim}
Le të jetë $f(x)$ dendësia synuar dhe $g(x)$ propozimi, me $f(x)\le Mg(x)$. Gjenerohet $X\sim g$ dhe pranohet nëse
\begin{equation}
 U\le\frac{f(X)}{Mg(X)}.
\end{equation}
Norma mesatare e pranimit është $1/M$. Propozimi duhet të mbulojë të gjithë mbështetjen e $f$ dhe bishtat e tij nuk duhet të jenë më të lehtë. Në dimension të lartë, refuzimi global bëhet shpesh joefikas.

\section{Kampionimi i rëndësisë}
Për një pritje nën $p$,
\begin{equation}
 \E_p[h(X)]=\int h(x)\frac{p(x)}{q(x)}q(x)\dd x
 \approx\frac1N\sum_{i=1}^Nh(X_i)w(X_i),\quad X_i\sim q.
\end{equation}
Propozimi $q$ duhet të vendosë mostra atje ku $|h|p$ është i madh. Peshat shumë të pabarabarta japin madhësi të vogël efektive të mostrës,
\begin{equation}
 N_{\mathrm{eff}}\approx\frac{(\sum_iw_i)^2}{\sum_iw_i^2}.
\end{equation}

\begin{kujdes}
Ndryshimi i farës nuk zëvendëson analizën statistike. Rezultati duhet të jetë i qëndrueshëm ndaj disa farave, ndërsa raportimi duhet të japë vlerësuesin dhe gabimin standard.
\end{kujdes}

\begin{lidhje}
Laboratori gjeneron uniforme, eksponenciale dhe Gauss; përdor histogramin, CDF empirike dhe autokorrelacionin; krahason transformimin e anasjelltë me refuzimin.
\end{lidhje}

\section*{Pyetje kontrolli}
\begin{enumerate}
\item Pse një histogram uniform nuk provon pavarësinë?
\item Çfarë ndodh nëse $Mg(x)<f(x)$ në një zonë?
\item Si organizohen vargjet pseudorastësore në llogaritje paralele?
\end{enumerate}
\section{Algoritmet e gjenerimit dhe kampionimit}
\subsection{Gjeneratori dhe gjendja e tij}
Në NumPy rekomandohet objekti \texttt{Generator}. Ai e bën burimin e rastësisë të dukshëm dhe shmang varësinë nga një gjendje globale.
\begin{lstlisting}[language=Python]
rng = np.random.default_rng(2026)
u = rng.random(1_000_000)
print(u.mean(), u.var())  # priten afersisht 1/2 dhe 1/12
\end{lstlisting}

\subsection{Transformimi i anasjelltë}
\begin{lstlisting}[language=Python,caption={Mostra eksponenciale pa përdorur funksionin e gatshëm.}]
def eksponenciale_anasjelltas(rng, lam, size):
    u = rng.random(size)
    return -np.log1p(-u) / lam

x = eksponenciale_anasjelltas(rng, lam=1.5, size=100_000)
print(x.mean(), 1/1.5)
\end{lstlisting}
Funksioni \texttt{log1p(-u)} është më i saktë kur $u$ është i vogël. Kontrollohen mesatarja, varianca dhe CDF-ja empirike.

\subsection{Box--Muller}
\begin{lstlisting}[language=Python,caption={Gjenerimi i çifteve Gauss standarde.}]
def box_muller(rng, n):
    u1 = rng.random(n)
    u2 = rng.random(n)
    r = np.sqrt(-2.0*np.log(u1))
    phi = 2.0*np.pi*u2
    return r*np.cos(phi), r*np.sin(phi)
\end{lstlisting}
Testohet jo vetëm histogrami i secilit komponent, por edhe kovarianca ndërmjet tyre.

\subsection{Kampionimi me refuzim}
\begin{lstlisting}[language=Python,caption={Refuzimi për një dendësi në interval të fundmë.}]
def refuzim(rng, f, a, b, M, n):
    pranuar = []
    prova = 0
    while len(pranuar) < n:
        x = rng.uniform(a, b)
        y = rng.uniform(0.0, M)
        prova += 1
        if y <= f(x):
            pranuar.append(x)
    return np.asarray(pranuar), n/prova
\end{lstlisting}
Këtu $M$ është kufi i sipërm i $f$ në $[a,b]$. Në formulimin me propozim të përgjithshëm, kontrollohet $f(x)\le M g(x)$.

\subsection{Testet minimale}
\begin{enumerate}
\item krahasohen momentet empirike me ato teorike;
\item vizatohet CDF-ja empirike;
\item testohet autokorrelacioni;
\item përsëritet me fara të ndryshme;
\item matet norma e pranimit dhe kostoja për mostër.
\end{enumerate}
Termi \emph{kampionim} përdoret për sampling dhe \emph{mostër} për sample. \emph{Gjenerator pseudorastësor} thekson natyrën deterministe të vargut.

Terminologjia probabilitare shqip mbështetet te \cite{capollari,butka,tasho}; algoritmet e gjenerimit dhe testimit krahasohen me trajtimet standarde \cite{press,robert}.

\subsection{Gjeneratorët pseudorastësorë si sisteme dinamike diskrete}
Një gjenerator pseudorastësor është një algoritëm determinist me gjendje të fundme. Gjeneratori kongruencial linear,
\begin{equation}
 x_{n+1}=(ax_n+c)\bmod m,
 \qquad u_n=\frac{x_n}{m},
\end{equation}
e bën këtë strukturë të dukshme. Për shkak se ka vetëm $m$ gjendje, vargu duhet të përsëritet; perioda dhe korrelacionet varen nga parametrat. Edhe kur histogrami i $u_n$ duket uniform, çiftet $(u_n,u_{n+1})$ mund të shtrihen në një numër të vogël hiperplanesh. Prandaj uniformiteti njëpërmasor nuk provon cilësinë shumëpërmasore.

Gjeneratorët modernë kanë gjendje shumë më të madhe dhe cilësi të testuar, por mbeten deterministë. Fara përcakton trajektoren. Ruajtja e farës siguron riprodhueshmëri, ndërsa përdorimi i së njëjtës farë për eksperimente që duhet të jenë të pavarura mund të krijojë korrelacione artificiale.

\subsection{Transformimi i anasjellë}
Nëse $U\sim\mathcal U(0,1)$ dhe $F$ është funksioni kumulativ i vazhdueshëm e rritës, atëherë
\begin{equation}
 X=F^{-1}(U)
\end{equation}
ka funksion kumulativ $F$, sepse
\begin{equation}
 P(X\le x)=P(F^{-1}(U)\le x)=P(U\le F(x))=F(x).
\end{equation}
Për eksponencialen $F(x)=1-e^{-\lambda x}$ merret
\begin{equation}
 X=-\frac{1}{\lambda}\ln(1-U).
\end{equation}
Shkrimi $\log1p(-U)$ është më i saktë kur $U$ është i vogël. Metoda është e thjeshtë kur inversi është i njohur; për shpërndarje të tjera nevojitet zgjidhje numerike ose metodë alternative.

\subsection{Box--Muller për mostrën Gauss}
Dy ndryshore Gauss standarde kanë dendësi të përbashkët
\begin{equation}
 f(z_1,z_2)=\frac{1}{2\pi}\exp[-(z_1^2+z_2^2)/2].
\end{equation}
Në koordinata polare, elementi i sipërfaqes është $r\dd r\dd\theta$, kështu që $\Theta$ është uniforme në $[0,2\pi)$ dhe $R^2$ ka shpërndarje eksponenciale me parametër $1/2$. Me $U_1,U_2$ uniforme të pavarura,
\begin{equation}
 R=\sqrt{-2\ln U_1},\qquad \Theta=2\pi U_2,
\end{equation}
dhe
\begin{equation}
 Z_1=R\cos\Theta,qquad Z_2=R\sin\Theta.
\end{equation}
Duhet shmangur $U_1=0$. Mostrat verifikohen jo vetëm me histogram, por edhe me mesatare, variancë, kuantile dhe korrelacion ndërmjet $Z_1$ e $Z_2$.

\subsection{Kampionimi me refuzim}
Le të jetë $f$ dendësia e synuar dhe $g$ dendësia propozuese, me $f(x)\le Mg(x)$. Gjenerohet $Y\sim g$ dhe $U\sim\mathcal U(0,1)$; $Y$ pranohet kur
\begin{equation}
 U\le\frac{f(Y)}{Mg(Y)}.
\end{equation}
Dendësia e një vlere të pranuar është proporcionale me
\begin{equation}
 g(y)\frac{f(y)}{Mg(y)}=\frac{f(y)}{M},
\end{equation}
pra pas normalizimit është $f$. Probabiliteti mesatar i pranimit është $1/M$. Një mbështjellëse e dobët e bën metodën të kushtueshme, sidomos në dimension të lartë.

\subsection{Kampionimi i rëndësisë dhe madhësia efektive}
Për një pritje nën $p$,
\begin{equation}
 I=\int h(x)p(x)\dd x
 =\int h(x)\frac{p(x)}{q(x)}q(x)\dd x,
\end{equation}
gjenerohen $X_i\sim q$ dhe përdoren peshat $w_i=p(X_i)/q(X_i)$. Kushti themelor është që $q>0$ kudo ku $|h|p$ kontribuon. Varianca mund të zvogëlohet shumë nëse $q$ vendos mostra në rajonet me kontribut të madh, por peshat shumë të pabarabarta e bëjnë vlerësimin të paqëndrueshëm. Një diagnostikë është
\begin{equation}
 N_{\mathrm{eff}}=\frac{(\sum_iw_i)^2}{\sum_iw_i^2}.
\end{equation}
Kjo nuk është identike me madhësinë efektive të një zinxhiri Markov, megjithëse të dyja matin humbjen e informacionit kundrejt $N$ mostrave ideale.

\subsection{Testimi i mostrave}
Një paketë minimale kontrollesh përfshin histogramin, CDF-në empirike, testin e momenteve dhe autokorrelacionin. Vlera-$p$ nuk provon rastësinë; ajo mat sa i pazakontë është statistiku nën një hipotezë të caktuar. Me shumë teste shfaqen rezultate të vogla edhe për një gjenerator të mirë, prandaj duhet parë modeli i përgjithshëm dhe jo vetëm minimumi i vlerave-$p$.

\subsection{Përmbledhje dhe ushtrime}
Gjenerimi i mostrës është pjesë e algoritmit shkencor. Një transformim i saktë matematikisht mund të dështojë numerikisht ose të trashëgojë korrelacionet e gjeneratorit bazë.
\begin{enumerate}
\item Gjeni kushtet për periodë të plotë të një gjeneratori kongruencial linear dhe testojini në parametra të vegjël.
\item Nxirrni transformimin e anasjellë për shpërndarjen Pareto.
\item Vërtetoni Box--Muller-in duke përfshirë Jakobinin polar.
\item Ndërtoni një mbështjellëse refuzimi për një dendësi trekëndore dhe parashikoni normën e pranimit.
\item Krahasoni dy propozime të kampionimit të rëndësisë me anë të variancës dhe $N_{\mathrm{eff}}$.
\end{enumerate}

\subsection{Rrymat e pavarura dhe ekzekutimi paralel}
Në një simulim paralel, përdorimi i farave $s,s+1,s+2,\ldots$ nuk garanton rryma të pavarura për çdo familje gjeneratorësh. Gjeneratorët modernë ofrojnë mekanizma për ndarje: kapërcim përpara, nënrryma ose ndërtim të gjendjeve fëmijë nga një sekuencë farash. Qëllimi është që asnjë punëtor të mos përdorë pjesë të mbivendosura të të njëjtit varg.

Rezultati duhet të jetë i riprodhueshëm edhe kur ndryshon numri i bërthamave. Një strategji është që çdo realizim të marrë farë të derivuar nga identifikuesi i realizimit, jo nga rendi i ekzekutimit. Atëherë planifikimi paralel mund të ndryshojë pa ndryshuar mostrat. Në raport ruhen fara bazë, metoda e ndarjes dhe versioni i bibliotekës.

Gjeneratori kriptografik nuk është domosdoshmërisht i nevojshëm për fizikë llogaritëse; kërkohen periodë e gjatë, cilësi statistikë, shpejtësi dhe ndarje e sigurt e rrymave. Për aplikime sigurie, parashikueshmëria deterministe është e papranueshme dhe kërkesat janë të tjera.

\subsubsection{Gjenerimi i vektorëve Gauss të korreluar}
Le të jetë $\vect Z\sim\mathcal N(0,I)$ dhe $\Sigma=LL^T$ faktorizimi Cholesky i kovariancës. Atëherë
\begin{equation}
 \vect X=\vect\mu+L\vect Z
\end{equation}
ka mesatare $\vect\mu$ dhe kovariancë
\begin{equation}
 \Cov(\vect X)=L\Cov(\vect Z)L^T=LL^T=\Sigma.
\end{equation}
Nëse $\Sigma$ është vetëm pozitivisht gjysmë e përcaktuar, Cholesky mund të dështojë; përdoret zbërthimi spektral dhe mbahen vlerat vetjake jo negative. Një vlerë vetjake zero përfaqëson kufizim të saktë linear dhe dimension efektiv më të vogël.

Kontrolli empirik kërkon shumë realizime dhe krahasim të mesatares e kovariancës së mostrës me objektivat. Devijimet priten të shkallëzohen me madhësinë e mostrës; nuk duhet kërkuar barazi e saktë. Grafikët e çifteve ndihmojnë të zbulohen shenja të gabuara ose renditje e gabuar e komponentëve.

\subsubsection{Metoda e përbërjes dhe shpërndarjet diskrete}
Nëse dendësia është përzierje
\begin{equation}
 f(x)=\sum_{k=1}^{K}w_kf_k(x),\qquad w_k\ge0,\quad\sum_kw_k=1,
\end{equation}
fillimisht gjenerohet komponenti $K$ me probabilitete $w_k$, pastaj $X\sim f_K$. Kjo metodë është e saktë dhe shpesh më efikase se refuzimi. Për shembull, një model zhurme me bërthamë të ngushtë dhe bisht të gjerë mund të simulohet si përzierje e dy shpërndarjeve Gauss.

Për shpërndarje diskrete me shumë kategori, kërkimi kumulativ kushton $\Oh(\log K)$ me kërkim binar. Metoda alias kryen parapërpunim $\Oh(K)$ dhe pastaj gjeneron çdo kategori në kohë $\Oh(1)$. Zgjedhja varet nga sa herë përdoret e njëjta shpërndarje.

\subsubsection{Efikasiteti i refuzimit në dimension të lartë}
Në njësi $d$-dimensionale, raporti i vëllimit të sferës ndaj kubit kufizues është
\begin{equation}
 \frac{V_d}{2^d}=\frac{\pi^{d/2}}{2^d\Gamma(d/2+1)},
\end{equation}
dhe shkon shpejt në zero. Prandaj gjenerimi uniform në kub dhe refuzimi jashtë sferës bëhet katastrofik në dimension të lartë. Ky shembull konkret tregon se një algoritëm intuitiv në dy përmasa mund të jetë i papërdorshëm në njëzet përmasa.

Transformimet që respektojnë gjeometrinë, kampionimi i rëndësisë ose MCMC shmangin këtë humbje. Megjithatë, MCMC prodhon korrelacion dhe kërkon diagnostikë; nuk ka metodë universale pa kosto.

\subsubsection{Testet si bateri, jo si vulë}
Një test uniformiteti kontrollon vetëm një aspekt. Testet seriale kontrollojnë çifte ose blloqe; testet spektrale kontrollojnë strukturën e rrjetës; testet e runs kontrollojnë vargjet rritëse dhe zbritëse. Kur kryhen shumë teste në nivel $5\%$, rreth $5\%$ pritet të dështojnë rastësisht. Duhet parë shpërndarja e vlerave-$p$, përsëritja në fara të ndryshme dhe rëndësia për algoritmin konkret.

Prova më bindëse për një simulim është kombinimi i një gjeneratori të njohur, testeve standarde dhe verifikimit të observablave ndaj rezultateve teorike. Një gjenerator mund të kalojë një bateri dhe përsëri të bashkëveprojë keq me strukturën periodike të një algoritmi të veçantë.

